Independent groups can maintain $\gscan c$ entries to find only $c$
winners. Sharing a safe upper bound lets them reduce this maintenance.
The bound is a complete key, including its tie position: once a group
holds $c$ distinct, fully scored candidates, its worst key bounds the
global cutoff. Other groups can reject keys later than that bound.
Scheduling separately chooses how many groups to run; both mechanisms
use a fixed refinement budget, chosen in \Cref{ch:refinement}.

\subsection{A complete-key upper bound with explicit witnesses}

Let $K_{(c)}$ be the global $c$th key when $|U|\geq c$. A group's
$c$ fully scored candidates are witnesses: their existence justifies an
upper bound. If their largest key is $\tau_j$, then
$K_{(c)}\preceq\tau_j$, because at least $c$ candidates precede or equal
it. Initialise the shared bound to the top sentinel and tighten it by
\begin{equation}
 \tau\leftarrow\min_{\prec}(\tau,\tau_j).
 \label{eq:5-shared-update}
\end{equation}
A group with fewer than $c$ valid candidates publishes no bound. If the
whole universe has fewer than $c$, all valid candidates remain eligible.
The next proposition guarantees that filtering by a complete published
bound preserves global winners, even if groups read older updates.

\begin{resultblock}
\begin{proposition}[Safe rejection by a witnessed bound]
\label{prop:5-shared-bound}
Assume fixed keys, publications backed by $c$ distinct candidates, and
reads that return either the initial sentinel or one complete published
key. Rejecting a candidate only when $K_q(x)\succ\tau$ cannot remove a
global top-$c$ candidate, irrespective of publication timing.
\end{proposition}
\begin{proof}
A publication has a set $W_\tau$ of $c$ distinct witnesses:
\[
 \max_{x\in W_\tau}K_q(x)\preceq\tau
 \quad\Longrightarrow\quad K_{(c)}\preceq\tau.
\]
Sentinel reads and minima of published keys obey the same inequality. Thus
\[
 x\in\operatorname{Top}_c(U)
 \ \Longrightarrow\ K_q(x)\preceq K_{(c)}\preceq\tau_{\rm seen}
 \ \Longrightarrow\ x\text{ is not rejected}.
\]
Stale reads only loosen the bound. Publications
made after earlier safe rejections remain valid: any $c$ retained
distinct candidates still supply $c$ witnesses, whether or not they are
the best $c$ of their group.
\end{proof}
\end{resultblock}
Combined with the coverage condition in \cref{prop:5-local-coverage},
this permits local and shared filtering together:
\begin{equation}
 K_q(x)\prec\lambda_j
 \quad\text{and}\quad
 K_q(x)\preceq\tau_{\rm seen},
 \label{eq:5-admission}
\end{equation}
Here $\lambda_j$ is the local compaction threshold and $\tau_{\rm seen}$
is the last shared key read. The local test is strict; the shared test
retains equality, including the candidate that supplied the bound.

\subsection{Publication, filtering and implementation limits}

The GPU encodes a complete key as one 64-bit integer. For nonnegative
FP32 scores, canonical positive zero and valid nonnegative positions,
\begin{equation}
 \kappa(D,p)=2^{32}\operatorname{uintbits}_{32}(D)+p.
 \label{eq:5-packed-key}
\end{equation}
On this non-NaN domain, comparing the packed integers has the same order
as comparing the score and then the position:
\[
 (D_1,p_1)\prec(D_2,p_2)\iff\kappa(D_1,p_1)<\kappa(D_2,p_2).
\]
Finite keys precede infinity and the all-one sentinel; invalid positions
cannot publish a bound.

\par\medskip\noindent\begin{minipage}{\linewidth}
\captionsetup{hypcap=false}
\centering
% One group loops over chunks; Q is shared across groups of one query.
\begin{tikzpicture}[font=\small\linespread{1}\selectfont,
 box/.style={draw=ThesisNavy!60,rounded corners=2pt,align=center,text width=5.1cm,minimum height=.85cm,inner sep=4pt},
 side/.style={box,text width=3.2cm,fill=ThesisNavy!5},
 arr/.style={-{Latex[length=2mm]},draw=ThesisNavy!80,thick}]
\node[box] (score) at (0,0) {Score next chunk\\Read shared bound once};
\node[box,fill=ThesisNavy!5] (admit) at (0,-1.5) {Key below local threshold\\and at or below shared bound?};
\node[box] (append) at (0,-3.0) {Append admitted keys};
\node[side] (discard) at (5.2,-1.5) {Discard candidate};
\node[box] (compact) at (0,-4.3) {If buffer needs space: compact\\Keep best $c$ keys};
\node[box] (publish) at (0,-5.8) {After compaction, with $c$ valid keys:\\publish worst retained key};
\node[box,fill=ThesisNavy!5] (more) at (0,-7.3) {More chunks?};
\node[side,text width=3.3cm] (shared) at (9.5,-4.0) {Shared query bound $Q$\\Initial value: no restriction\\Update: atomic minimum};
\node[box,text width=6.1cm] (final) at (7.5,-7.3) {Final compact; publish if $c$ valid\\Trim to latest shared bound\\Send retained prefix to merge};
\draw[arr] (score)--(admit);\draw[arr] (admit)--node[right] {yes}(append);
\draw[arr] (admit)--node[above] {no}(discard);
\draw[arr] (append)--(compact);
\draw[arr] (discard.south)--(5.2,-3.65)--(0,-3.65);
\draw[arr] (compact)--(publish);\draw[arr] (publish)--(more);
\draw[arr] (more.west)--node[above,pos=.45] {yes}(-3.3,-7.3)--(-3.3,0)--(score.west);
\draw[arr] (more.east)--node[above] {no}(final.west);
\draw[arr] (publish.east)--(9.5,-5.8)--(shared.south);
\draw[arr,dashed] (shared.north)--(9.5,0)--(score.east);
\draw[arr] (final.east)--(11.6,-7.3)--(11.6,-4.0)--(shared.east);
\end{tikzpicture}

\captionof{figure}[Shared-key filtering over scan chunks]{Shared-key filtering over scan chunks.}
\label{fig:5-shared-bound}
\end{minipage}\par\medskip
A designated thread reads the bound once per scanned chunk. The block
first computes full primary scores, then admits eligible keys and
compacts when space is needed. With $c$ valid keys after compaction,
an atomic minimum can publish its worst retained key. After the last
chunk, final compaction and trimming against the latest bound produce
the prefix sent to merging. An older read may leave extra candidates,
which exact merging removes. Reset, scan and merge are ordered in one
GPU stream; groups do not wait for one another.

The proof requires an aligned, whole-key read. Generated 64-bit loads
and atomic updates support this premise under the PTX memory model
\cite{ptx_isa}; C++ \texttt{volatile} alone does not provide it
\cite{cuda_programming_guide}.\footnote{The check uses CUDA 13.0.88
\texttt{ptxas} for \texttt{sm\_120}, with
\path{ld.volatile.global.u64} and \path{atom.global.min.u64}. It did not
capture the native instructions produced by the driver's runtime JIT.}
As \cref{fig:5-shared-bound} shows, filtering follows full primary
scoring. It saves buffer and merge maintenance while adding bound reads
and atomic updates; \cref{sec:7-4} measures this trade-off.

\subsection{Choosing parallelism without changing the budget}

Using $b\gscan$ blocks for a batch of $b$ queries exposes more parallel
work, but each group also needs a table, buffer and merge input. The
scheduler balances extra parallelism against this repeated work. Define
the block size, device thread capacity and variation in candidate counts:
\[
 T_{\rm blk}=\text{scan threads/block},\quad
 H=n_{\rm SM}T_{\rm resident/SM},\quad
 v=\frac{\operatorname{sd}(n_c)}{\operatorname{mean}(n_c)}.
\]
The fill rule adds groups when the batch would otherwise use only part
of that capacity. A second rule adds groups when query workloads are
uneven:
\begin{align}
 g_{\rm fill}
 &=\max\!\left(\{1\}\cup
       \{h\in\{2,4,8\}:bhT_{\rm blk}\leq0.6H\}\right),
       \label{eq:5-fill-rule}\\
 g_{\rm skew}&=
   \begin{cases}4,&v\geq0.3,\\1,&v<0.3,\end{cases}
 \qquad
 g_{\rm chosen}=\max(g_{\rm fill},g_{\rm skew}).
 \label{eq:5-skew-rule}
\end{align}
The skew rule can override the fill limit. To estimate $v$, the evaluation
uses 512 database vectors as pseudoqueries under fixed routing. It needs
no GT and is separate from Chapter 6's output-agreement calibration.
Later searches use the stored statistic and batch size without an extra
host readback. The constants were fitted on four pilot configurations,
then tested on further settings. The one-time scheduler calibration is
outside steady-state timing; its recorded setup times are in
\cref{sec:a-supplement}.

Choosing one group dispatches $\gscan=0$ with the same index, probes and
budget. Its cutoff-tie behaviour is described in \cref{sec:5-4}; integrated
gains and scheduling losses are evaluated in \cref{sec:7-6}.

\subsection{Observed execution paths}

The batch size explains which path supplies the main cross-system results.
For 1,000 queries with 512 scan threads, two groups already exceed the
fill target on the 170-SM RTX 5090, which supports 48 warps per SM
\cite{blackwell_tuning}:
\[
 1000\cdot2\cdot512=1{,}024{,}000
 >0.6\cdot170\cdot48\cdot32=156{,}672.
\]
The fill rule therefore chooses one group, while the skew rule can still
choose four. In the main comparison, all 9 of 38 routing records above $v=0.3$ are Vogue
($0.338$--$0.622$). Its 27 settings use four groups; all 105 arXiv/OpenAI
settings use $\gscan=0$. Those settings benefit from LUTs, packing,
cooperative refinement and probe selection, rather than grouped merging.
Vogue comparisons show 3--15 cutoff-tie differences per 1,000 queries,
no non-tie differences and identical final outputs; at 128 probes the
tie range is 7--11.

The budget study uses 128-query calibration samples and larger held-out
batches. Their path counts are
\[
 \begin{array}{c|rrr}
 &G=0&G=2&G=4\\\hline
 \text{held-out}&10&0&2\\
 S=128&2&8&2
 \end{array}
\]
Eight of twelve settings switch paths. The benefit depends on both
batch size and workload skew, rather than a fixed batch-256 boundary.
The heuristic does not model residency, compaction or merge cost, and
pseudoquery statistics can become outdated. Its regret is the speedup
of the best tested fixed group count divided by the scheduler's speedup;
the worst observed value is 1.217 at batch 256. Refinement budgets remain
fixed in these tests; \cref{sec:5-6} describes the refinement itself.
